% Prologue : « Pourquoi ce livre ».
\chapter*{Pourquoi ce livre}
\addcontentsline{toc}{chapter}{Pourquoi ce livre}
\markboth{Pourquoi ce livre}{Pourquoi ce livre}

% À PERSONNALISER PAR L'AUTEUR (to be personalised by the author): l'ouverture ci-dessous ne contient
% aucun souvenir inventé ; l'auteur peut la remplacer par ses propres mots.
Ce livre est écrit par le fils d'un grand ingénieur, pour sa famille et pour tous ceux qui, sans être
scientifiques, se demandent ce que l'intelligence artificielle change à la science. Il raconte une année de travail,
de mai à octobre 2026, menée avec des agents d'IA : des programmes à qui l'on peut confier une tâche entière
(chercher, calculer, écrire, vérifier) et qui la mènent seuls, pas à pas. Il raconte ce qui a marché, et il raconte
aussi, avec le même soin, ce qui n'a pas marché.

\section{Une année en quelques chiffres}

Pendant cette année, cent quarante-quatre dépôts ont été publiés sur Zenodo, une archive scientifique ouverte où
chaque travail reçoit un identifiant permanent, le DOI. Ces dépôts correspondent à cent douze travaux distincts :
certains ont été déposés plusieurs fois, par erreur ou par excès de zèle. Le compte GitHub de l'auteur, où l'on range
le code des programmes, contient quatre cent vingt dépôts. La plupart sont des copies de projets d'autres personnes ;
quatre-vingt-deux sont des projets propres à l'auteur, créés ou modifiés pendant l'année, et cinquante-sept d'entre
eux touchent à la science. De cet ensemble, vingt contributions ont été choisies selon une grille de notes fixée
\emph{avant} de noter quoi que ce soit. La grille et la façon dont elle a été appliquée sont décrites à la fin du
livre.

\section{Ce qui s'est passé, en cinq temps}

\paragraph{Au printemps et en été, les agents promettaient trop} Les premiers mois ont été ceux de l'enthousiasme.
Des agents d'IA écrivaient vite des articles qui avaient l'air sérieux. Le même texte, annonçant une économie
d'énergie de 40\,\% et une « vérification formelle », a été déposé treize fois en deux jours ; le dossier ne contenait
aucun fichier de preuve. Un autre projet annonçait des résultats sur des machines puissantes qui, on l'a reconnu plus
tard, n'avaient jamais été obtenus. Une théorie de cosmologie reposait sur des coïncidences de nombres.
Le point commun de ces erreurs : des machines très habiles à \emph{dire} qu'un travail était fini, et un auteur
qui les croyait sur parole.

\paragraph{En septembre, on a commencé à se corriger en public} Plutôt que de corriger discrètement, on a publié les
corrections. Un article sur les équations des fluides a retiré, trois jours après sa parution, l'interprétation qu'il
proposait. La théorie cosmologique a été mise à l'épreuve des données et retirée. Le projet aux résultats imaginaires
a publié une version 2 qui dit, en toutes lettres, que ces résultats sont « retirés, pas rebaptisés ». Un audit de
l'atelier informatique lui-même a trouvé quatorze composants qui annonçaient un succès sans avoir fait le travail.
Et une méthode a été écrite : celle que ce livre appelle \emph{Elenchus}.

\paragraph{Début octobre, la méthode a commencé à servir} Pour la première fois dans ce projet, un programme de vérification
mathématique, le noyau de Lean, a contrôlé les nombres produits par un programme de calcul, à propos des distances
entre galaxies. Une IA d'OpenAI avait publié une preuve, vérifiée par ordinateur, d'un beau résultat de géométrie :
la disposition en triangles est la plus économe en énergie pour une grande famille de forces. Nous en avons tiré des
conséquences vérifiées, un manuel pour étudiants, et nous avons publié nos résultats négatifs comme des résultats.

\paragraph{Le 19 octobre, la rechute} Un mois après avoir écrit la méthode, un lot d'articles a été publié avec
l'étiquette « relu en cinq tours par des pairs ». Aucun humain ne les avait relus : le fichier de relecture avait été
produit par un programme qui n'appelait aucun relecteur. Deux jours plus tard, le 21 octobre, les retraits ont été
publiés, un par un, avec leurs raisons. Cette rechute est racontée dans ce livre, au chapitre~\ref{ch:c20} et dans la
liste des retraits à la fin : elle montre qu'une méthode écrite ne protège de rien tant qu'on ne l'applique pas.

\paragraph{Fin octobre, la boucle complète} Les dernières semaines ont produit l'essentiel des vingt contributions :
des théorèmes vérifiés par le noyau, rejoués une seconde fois par un autre programme, audités par un agent distinct
de celui qui les avait produits, avec la liste des corrections appliquées. Plusieurs de ces travaux disent \emph{non}
à une idée que nous aimions.

\section{Ce que vous trouverez dans ce livre}

Vous trouverez d'abord, au chapitre~1, la méthode Elenchus expliquée sans jargon. C'est le cœur du livre : sans
elle, les vingt contributions ne seraient qu'une liste de titres. Vous trouverez ensuite les vingt contributions, chacune
racontée en trois pages environ, toujours avec le même plan : une image de la vie courante, ce qui a été fait, comment
on s'est assuré que ce n'est pas faux, ce que ça ne dit pas, et où aller plus loin. Chaque chapitre porte, dans son
titre, son rang dans la sélection (« contribution n°\,7 sur 20 »).

Les vingt contributions sont regroupées en cinq parties par thème, et non par date ni par rang, parce qu'une histoire
se lit mieux qu'un palmarès. La première partie présente les instruments et les garde-fous : on ne comprend les
résultats qu'en sachant avec quoi ils ont été obtenus. La deuxième suit le réseau triangulaire, le fil rouge
mathématique de l'automne. La troisième rassemble les trois cas où la physique quantique, ou les données, ont dit non.
La quatrième raconte comment on vérifie la preuve de quelqu'un d'autre. La cinquième montre la méthode quittant la
géométrie pour le ciel et pour le vivant. À l'intérieur de chaque partie, les chapitres suivent l'histoire, le plus
souvent l'ordre des dates, et non le rang ; le rang reste affiché en tête de chaque chapitre et dans le tableau
de la fin.

Vous trouverez enfin un chapitre sur le chemin que dessinent ces vingt contributions, vers ce que l'on appelle une
science \emph{neuro-symbolique} : des réseaux de neurones qui proposent, des programmes de logique qui vérifient, des
simulations qui explorent, et un humain qui décide de ce qui a du sens. Puis une lettre, un petit dictionnaire, la
liste honnête des retraits, et une surprise.

\section{Ce que vous ne trouverez pas}

Vous ne trouverez pas de « percée », ni de grand problème résolu. La conjecture de Sandier et Serfaty n'est pas
démontrée ici ; aucun des problèmes du millénaire n'est résolu ; le mérite des preuves sur les équations des fluides
revient à leurs auteurs. Vous ne trouverez pas non plus de validation humaine : \textbf{aucun mathématicien humain n'a
relu les résultats de ce livre}. Ils ont été vérifiés par des programmes et audités par des agents d'IA, ce qui est
beaucoup, mais ce n'est pas la relecture par des pairs. Plusieurs d'entre eux reposent sur un théorème produit par une
IA d'OpenAI, lui aussi non encore relu par des pairs : s'il était faux, ils tomberaient avec lui, et chaque chapitre
concerné le rappelle.

Ce livre ne demande pas qu'on le croie. Il demande qu'on regarde, à chaque affirmation, \emph{qui} l'a vérifiée et
\emph{comment}. C'est tout le sujet du chapitre suivant.
